\DocumentMetadata{
  pdfversion=2.0,
  pdfstandard=ua-2,
  lang=en-US,
  tagging=on,
  tagging-setup={math/setup=mathml-SE,tabsorder=structure}
}
\documentclass[11pt,letterpaper]{article}
\usepackage[margin=0.68in,includefoot]{geometry}
\usepackage{newcomputermodern}
\usepackage{amsmath,amscd,mathtools}
\providecommand{\square}{\mdlgwhtsquare}
\usepackage[hidelinks]{hyperref}
\hypersetup{
  pdftitle={Algebra First-Year Exam, May 2023, Part 2},
  pdfauthor={University of Florida Department of Mathematics},
  pdfsubject={Graduate mathematics examination},
  pdfdisplaydoctitle=true,
  bookmarksopen=true
}
\setcounter{secnumdepth}{0}
\setlength{\parindent}{0pt}
\setlength{\parskip}{0.28em}
\setlength{\emergencystretch}{3em}
\setlength{\leftmargini}{2.15em}
\setlength{\leftmarginii}{2.65em}
\setlength{\leftmarginiii}{3em}
\renewcommand{\labelenumi}{\textbf{\theenumi.}}
\pagestyle{empty}
\raggedbottom
\allowdisplaybreaks
\newenvironment{examproblems}[1]{%
  \begin{enumerate}%
  \setcounter{enumi}{#1}%
  \setlength{\topsep}{0.35em}%
  \setlength{\partopsep}{0pt}%
  \setlength{\itemsep}{0.48em}%
  \setlength{\parsep}{0.18em}%
}{\end{enumerate}}
\newenvironment{examsubparts}{%
  \begin{enumerate}%
  \setlength{\topsep}{0.18em}%
  \setlength{\partopsep}{0pt}%
  \setlength{\itemsep}{0.16em}%
  \setlength{\parsep}{0.1em}%
}{\end{enumerate}}
\newenvironment{examinstructions}{%
  \par\begingroup\itshape
}{%
  \par\endgroup\vspace{0.18em}
}
\makeatletter
\renewcommand\section{\@startsection{section}{1}{0pt}%
  {-1.25ex plus -.25ex minus -.1ex}%
  {0.55ex plus .1ex}%
  {\normalfont\large\bfseries}}
\renewcommand{\@maketitle}{%
  \newpage\null\vskip -1.2em%
  {\footnotesize\bfseries
    \href{https://www.ufl.edu/}{University of Florida}, \href{https://math.ufl.edu/}{Department of Mathematics}\par}%
  \vskip 0.18em%
  \tagstructbegin{tag=Title}%
    {\LARGE\bfseries\href{https://gma.math.ufl.edu/exams/algebra-first-year/}{Algebra First-Year Exam}, May 2023, Part 2\par}%
  \tagstructend%
  \vskip 0.35em%
  \tagmcbegin{artifact}%
    \hrule height 1pt%
  \tagmcend%
  \vskip 0.55em%
}
\makeatother
\title{Algebra First-Year Exam, May 2023, Part 2}
\author{}
\date{}
\begin{document}
\maketitle
\thispagestyle{empty}
\begin{examinstructions}
Answer four problems. If you turn in more than four, only the first four will be graded. Within reason, you may use theorems as long as you state them clearly.
\end{examinstructions}
\begin{examproblems}{0}
\item
Show that the ring 
\[
\mathbb{Z}[\sqrt{-5}]=\{a+b\sqrt{-5}\mid a,b\in\mathbb{Z}\}
\]
 is not a unique factorization domain.
\item
Which of the following rings are integral domains? Fields? Explain.
\begin{examsubparts}
\item[\textnormal{(i)}]
\(\mathbb{Q}[X,Y]/(X(X+1)+2X^2Y-Y^3)\)
\item[\textnormal{(ii)}]
\(\mathbb{Q}[X,Y]/(X^2-Y^2)\)
\item[\textnormal{(iii)}]
\(\mathbb{Z}[X,Y]/(2,X,Y-1)\).
\end{examsubparts}
\item
Let~\(R\) be an integral domain. Recall that an~\(R\)-module~\(M\) has rank~\(n\) if~\(n\) is the largest integer such that~\(M\) has~\(n\) \(R\)-linearly independent elements. Show that an~\(R\)-module~\(M\) has rank~\(n\) if and only if~\(M\) has a submodule \(N\subseteq M\) such that~\(N\) is free of rank~\(n\), and \(M/N\) is torsion.
\item
Let~\(F\) be a field, \(M\) an extension of~\(F\) and \(L,K\) two finite extensions of~\(F\) contained in~\(M\). Show that if \([L:F]\) and \([K:F]\) are relatively prime, \(L\cap K=F\).
\item
This problem has two parts.
\begin{examsubparts}
\item[\textnormal{(i)}]
List all irreducible polynomials in \(\mathbb{F}_2[X]\) of degree at most three.
\item[\textnormal{(ii)}]
List representatives for each conjugacy class of the group \(GL_3(\mathbb{F}_2)\).
\end{examsubparts}
\item
Let~\(F\) be a field, \(V\) an~\(F\)-vector space of finite dimension and \(T:V\to V\) a linear transformation. Show that if the minimal polynomial of~\(T\) is irreducible of degree~\(d\) then the dimension of~\(V\) is a multiple of~\(d\).
\end{examproblems}
\end{document}
